\paragraph{Removing a hyperbolic plane on the given lattice.}
The integral part of the change of model admits the following direct
construction. It will allow the local descent to use a rational plane
while retaining the lattice at every finite place.

\begin{proposition}\label{famtheta:plane}
Let \(E/F\) be the CM extension in the theta tower, unramified at
every finite place. Let \(L\) be a lattice in an \(E\)-space of
dimension \(n\ge3\), with nondegenerate Hermitian or skew-Hermitian
form \(b\), linear in its first argument. Suppose \(L\) is self-dual
for an invertible conjugation-stable fractional ideal \(\mathfrak d\)
and the form has both signs at each infinite place (after multiplication
by an imaginary scalar in the skew case). There are rational isotropic
vectors \(x,y\) with \(b(x,y)=1\) for which
\[
 L=(L\cap Ex)\ \oplus\ (L\cap Ey)\ \perp\ L_0 .
\]
Here the first two lines are paired perfectly with values in
\(\mathfrak d\), and \(L_0\) is self-dual for \(\mathfrak d\).
At any prescribed finite set of split places where
\(\mathfrak d_w=\mathcal O_{E,w}\) on both sides, the two vectors
may simultaneously be taken as local generators of their lines.
\end{proposition}
\begin{proof}
In the skew case multiply \(b\) temporarily by a nonzero element
negated by conjugation. The associated \(F\)-quadratic form
\(v\mapsto b(v,v)\) for the resulting Hermitian space is nondegenerate
of dimension \(2n\). It is indefinite at every real place and isotropic
at every finite place, since quadratic forms over a finite extension
of \(\mathbf Q_p\) of dimension at least five are isotropic (also for
\(p=2\)). Hasse--Minkowski supplies a rational isotropic \(E\)-line
\(U\) for the original form. The intersection \(I=L\cap U\) is
saturated and of rank one. Both \(I\) and \(L/I\) are torsion-free
finite modules over the Dedekind ring \(\mathcal O_E\), hence projective.
In particular \(I\) is a direct summand as a module.

Put \(J=\operatorname{Hom}_{\mathcal O_E}(\bar I,\mathfrak d)\),
where the bar denotes the conjugate module. Self-duality identifies
\(L\) with \(\operatorname{Hom}(\bar L,\mathfrak d)\) by the pairing.
Restriction to \(\bar I\) is onto by the direct summand property.
Choose an integral section \(s:J\to L\) of this restriction,
using projectivity of \(J\). The pairing of \(sJ\) and \(I\)
is perfect in both directions, by biduality of the invertible modules
and conjugation stability of \(\mathfrak d\).
Choose \(\theta\in\mathcal O_E\) of trace one. Indeed the trace ideal
is the unit ideal locally at every finite prime: at a split prime
this is addition, and in an unramified quadratic extension it follows
by reduction from the surjective residue-field trace. This argument
also covers residue characteristic two.
There is a unique linear map \(C:J\to I\) with
\[
 b(Cj,sj')=\theta\, b(sj,sj')\qquad (j,j'\in J)
\]
by the perfect pairing. Hermitian symmetry with its given sign implies
\(b(sj,Cj')=\bar\theta b(sj,sj')\). Since \(b(I,I)=0\), the map
\(s'=s-C\) is another section and \(b(s'J,s'J)=0\).
The Gram pairing on \(I\oplus s'J\) has zero diagonal blocks and
perfect off-diagonal blocks. It is self-dual for \(\mathfrak d\).
For any \(z\in L\), its pairing functional on this submodule is thus
realized by a unique vector inside it. Subtracting that vector gives
the orthogonal projection to \(L_0\). This proves the integral
decomposition and self-duality of the complement. The two summands
of the plane are precisely the full lattice intersections with their
rational spans.

Choose a rational generator \(x\) of \(I\otimes E\); by weak
approximation scale it to generate \(I\) at all the indicated primes.
There is a unique \(y\in(s'J)\otimes E\) with \(b(x,y)=1\). Perfection
and \(\mathfrak d_w=\mathcal O_{E,w}\) make it a generator at the
same places, including the conjugate primes. Complete these two vectors
by bases of \(L_0\) at the split places and put them in the designated
columns. At each real place the plane has signature \((1,1)\);
choosing a Cartan involution on it and on its orthogonal complement
gives the adapted marking. The lower signatures lose one of each sign.
\end{proof}
\src{src/Development/WeilBound/ThetaLifts/TowerModels/PlaneMain.lean}{47}{rational isotropy and integral trace one},
\src{src/Development/WeilBound/ThetaLifts/TowerModels/PlaneSkew.lean}{18}{skew form and codifferent}, and
\src{src/Development/WeilBound/ThetaLifts/TowerModels/TowerCodimensionTwoOf.lean}{32}{marked planes for the two towers}.
